\section{Introduction}

The capacity of the hippocampus to support vivid recollection of individual experiences is among the most striking features of mammalian memory \cite{scoville1957loss}. Within the hippocampal formation, the subfield CA3 has long been regarded as an autoassociative network in which dense recurrent excitation among pyramidal cells, together with spike-timing-dependent plasticity at those synapses \cite{bi1998synaptic}, permits the reinstatement of stored activity patterns from noisy or partial cues \cite{mcnaughton1987hippocampal,hopfield1982neural}. Sensory information destined for CA3 enters this circuit along two anatomically segregated routes from the entorhinal cortex (EC). A direct perforant path (PP) projects from EC layer II onto the distal dendrites of CA3 pyramidal cells, while an indirect route passes first through the dentate gyrus (DG) before its mossy fibers (MF) contact CA3 on more proximal dendrites. The functional justification for this dual-input architecture remains unsettled.

The two pathways deliver information with systematically different statistical properties. DG is equipped with a large population of granule cells and a high tonic inhibitory tone, which keeps their activity sparse and interference among overlapping representations low \cite{engin2015tonic,leutgeb2007pattern}. Sparsification in feedforward projections is in turn associated with decorrelation, so that MF inputs to CA3 inherit both a sparse and a decorrelated character. PP connectivity, by contrast, is dense, and natural correlations between similar sensory inputs should therefore survive the projection to CA3. Previous theories of the circuit have proposed that the MF pathway teaches CA3 to store well-separated memory traces during encoding, while the PP pathway mediates retrieval from a single hybrid attractor \cite{treves1992computational}. Under these accounts, each stored experience is represented by one pattern at equilibrium, irrespective of which route was used to deliver it.

Here we consider an alternative possibility. We ask whether CA3 can store and retrieve two distinct encodings of the same memory---one MF-like and one PP-like---and whether the physiological theta oscillation, which rhythmically modulates inhibitory tone across phases \cite{tsodyks1996population}, serves as the switch that selects between them. If this were true, sparse, decorrelated MF encodings and dense, correlated PP encodings would not collapse into a single stored pattern; instead, the network would multiplex the two, with different inhibitory thresholds favoring different retrieval modes. Conceptually, MF attractor basins would remain separate and support pattern separation suitable for distinguishing examples, while PP basins would merge along shared features and support concept-like generalization across similar experiences \cite{quiroga2005invariant,knowlton1993learning,zeithamova2019brain}.

This paper develops the proposal in three stages. First, we construct a biologically motivated model in which EC, DG, MF and PP encodings are produced by feedforward networks with realistic statistics, and in which a Hopfield-like CA3 network \cite{hopfield1982neural,tsodyks1988enhanced} stores linear combinations of MF and PP patterns. Asynchronous retrieval dynamics controlled by a threshold \cite{amit1985spin} are shown to alternate the network between distinct MF-like example representations at high threshold and merged PP-like concept representations at low threshold. Second, we test the core prediction of this account, that single CA3 place cells should convey more information per spike about fine, example-like positions during sparse theta phases, and more information about coarse, concept-like positions during dense phases, using publicly available rodent electrophysiological recordings \cite{mizuseki2013hc3,karlsson2015hc6}. Third, we abstract the principle to artificial neural networks: a DeCorr loss term that decorrelates batch representations, and a HalfCorr variant that applies the decorrelation to only half of the final hidden layer, are introduced and shown to govern a tradeoff between concept-like and example-like learning in a multilayer perceptron trained on MNIST \cite{lecun1998gradient}.

The contributions of this work are: (i) an autoassociative model in which a single CA3 network stores both MF and PP encodings of each memory and alternates between them under threshold control, together with a closed-form relationship between presynaptic and postsynaptic statistics for the feedforward pathways; (ii) empirical confirmation of the model's signature predictions in CA3 place fields across two datasets and two behavioral tasks, together with a Bayesian population analysis showing that sparse theta phases support more confident and more accurate decoding of behaviorally relevant variables; and (iii) a plug-and-play loss function that imports the same computational principle into artificial networks and improves multitask performance on digit classification and set identification.

\section{Related Work}

\paragraph{Autoassociative memory and the hippocampus.} The idea that CA3 performs pattern completion through recurrent excitation traces back to foundational work in attractor networks \cite{hopfield1982neural,tsodyks1988enhanced,amit1985spin} and to the complementary learning systems framework \cite{mcnaughton1987hippocampal,mcclelland1995why}. Several computational accounts have specifically addressed the role of the MF pathway, typically concluding that MF inputs support the storage of well-separated patterns that are later retrieved through PP cues \cite{treves1992computational}. Our model differs in asserting that two distinct encodings per memory coexist in CA3 and are retrieved in alternation, rather than collapsing into a single attractor.

\paragraph{Pattern separation in the dentate gyrus.} Experimental work has confirmed that DG decorrelates inputs to CA3 more aggressively than CA3 does on its own, by recruiting largely non-overlapping granule cell populations when environments differ \cite{leutgeb2007pattern} and by maintaining the high tonic inhibition that keeps activity sparse \cite{engin2015tonic}. These findings motivate our choice of densities (aMF = 0.02, aPP = 0.2) and within-concept correlations (rhoMF = 0.01, rhoPP = 0.09) for MF and PP encodings in the model.

\paragraph{Theta phase coding and place cells.} Phase precession, in which spikes advance to earlier theta phases as an animal traverses a place field, is a robust property of hippocampal place cells \cite{okeefe1993phase,tsodyks1996population} and is accompanied by phase-dependent changes in tuning sharpness \cite{skaggs1996theta}. Our analysis extends this literature by partitioning theta into sparse and dense halves and asking, separately for fine and coarse position binning, which half of the cycle conveys more sparsity-corrected information. The prediction tested is the unique signature of a threshold-switched multiplexed code.

\paragraph{Concepts, categories, and the medial temporal lobe.} Invariant responses of single human medial temporal lobe neurons to different depictions of a particular concept \cite{quiroga2005invariant}, together with behavioral and imaging evidence for hippocampal involvement in categorization \cite{knowlton1993learning,zeithamova2019brain}, establish that the hippocampus does more than barcode experiences. The prototype- and clustering-based accounts of category learning \cite{anderson1991adaptive,love2004sustain} provide a natural target representation for our PP attractors, which merge similar examples along shared features.

\paragraph{Representation shaping in artificial neural networks.} Regularizers that shape hidden representations have been used primarily to reduce overfitting. The DeCov loss, for example, decorrelates pairs of neurons across all inputs \cite{cogswell2016reducing}, whereas our DeCorr loss decorrelates pairs of inputs across all neurons in a layer. The two are complementary rather than equivalent: DeCov improves generalization on the concept task, while DeCorr induces the opposite effect by encouraging representations that discriminate examples. Our HalfCorr variant combines correlated and decorrelated components within a single layer to resolve the resulting tradeoff, bridging a biological motivation drawn from CA3 \cite{hopfield1982neural,schapiro2017complementary} with a concrete machine-learning construction applicable to networks trained with standard optimizers \cite{kingma2015adam,lecun2015deep}.

\paragraph{Cortical inhibition as threshold control.} Our interpretation of theta as a subtractive, phase-varying threshold echoes the broader view that inhibitory interneurons shape cortical activity through both subtractive and divisive effects \cite{isaacson2011inhibition}. The specific prediction tested here---that the sparser, more inhibited half of the cycle carries sharper tuning to specific features, while the denser half carries broader, concept-like tuning---is distinctive of a subtractive mechanism and is the central experimental signature of the proposed multiplexed code.
